% Compile with latexmk (uses ~/.latexmkrc: LuaLaTeX, output in build/)
% Week 9: probability, conditional probability, Bayes' formula, odds,
% independence, discrete random variables, expectation and variance.
% Sessions 1 and 2 are held on the SAME DAY, back to back (2 x 90 min,
% short break in between).
% Reference: Ross, A First Course in Probability, 8th ed., Ch. 2--4 (notation
% and the problems paraphrased); the course lecture slides only for comparison.
% The exercises are in exercises/ (one file each), placed by \exerciseframe.

\documentclass[aspectratio=169,11pt,usepdftitle=false]{beamer}
\usepackage{fontspec}
\setsansfont{Latin Modern Sans}
\usepackage{amsmath,amssymb,mathtools}
\usepackage{tikz}
\usetikzlibrary{calc}
\usepackage{booktabs}
\usepackage{cancel}

\usetheme{Madrid}

\definecolor{navy}{RGB}{24,48,75}
\definecolor{teal}{RGB}{0,120,125}
\definecolor{gold}{RGB}{190,130,20}

\setbeamercolor{structure}{fg=navy}
\setbeamercolor{alerted text}{fg=teal}
\setbeamercolor{block title}{bg=navy,fg=white}
\setbeamercolor{block body}{bg=navy!5}
\setbeamercolor{block title example}{bg=teal,fg=white}
\setbeamercolor{block body example}{bg=teal!6}
\setbeamercolor{block title alerted}{bg=gold,fg=white}
\setbeamercolor{block body alerted}{bg=gold!8}

\setbeamertemplate{navigation symbols}{}
\setbeamertemplate{footline}{%
  \usebeamerfont{footline}\hspace{1em}\insertshortauthor\hfill%
  \insertshorttitle\quad|\quad%
  \insertframenumber/\inserttotalframenumber%
  \hspace{1em}\vspace{2mm}%
}
\setbeamertemplate{itemize items}[circle]

% Contact (beamer loads hyperref itself)
\newcommand{\myemail}{tiborsandor2001@gmail.com}

\newcommand{\N}{\mathbb{N}}
\newcommand{\Om}{\Omega}
% Expectation, variance and standard deviation, as in Ross: E[X],
% Var(X), SD(X)
\newcommand{\E}{\operatorname{E}}
\newcommand{\Var}{\operatorname{Var}}
\newcommand{\SD}{\operatorname{SD}}
\newcommand{\sol}{\alert{Solution.}\ }

% Small Venn diagram of two events E, F in Omega; #1 = shading commands
\newcommand{\venn}[1]{%
  \begin{tikzpicture}[scale=.42,baseline=(O.base)]
    \node (O) at (.5,0) {\phantom{x}};
    #1
    \draw[thick] (-1.5,-1.15) rectangle (2.5,1.15);
    \draw[thick,navy] (0,0) circle (.85);
    \draw[thick,navy] (1,0) circle (.85);
    \node at (-.35,0) {\scriptsize $E$};
    \node at (1.35,0) {\scriptsize $F$};
    \node at (2.15,.8) {\scriptsize $\Omega$};
  \end{tikzpicture}}
\newcommand{\shadeall}[1]{\fill[#1] (-1.5,-1.15) rectangle (2.5,1.15);}

% The exercises (one file each in exercises/) and \exerciseframe
% Exercises of Week 9, shared by beamer.tex, the study guide, the script and
% exercises.tex. Every exercise lives in a file of its own,
% exercises/NN-slug.tex, numbered 9.NN (see exercises.lua for the docstring).
% The body of an exercise file:
%
%   \begin{problem} ... \end{problem}      the statement, used everywhere;
%                                          inside it, \begin{context} ...
%                                          \end{context} is background that
%                                          the slides leave out
%   \begin{slideproblem} ... \end{slideproblem}
%                                          (optional) a shorter statement
%                                          for the slides
%   \begin{slides} ... \end{slides}        one frame of the deck: \exstatement
%                                          places the statement; repeat the
%                                          environment for more frames
%   \begin{solution} ... \end{solution}    the detailed solution (study
%                                          guide, exercise PDFs, website)
%
% Commands for the documents (the argument is the slug: 01-license-plates.tex
% has the slug license-plates):
%
%   \exref{slug}                 the number, e.g. 9.1, for "Exercise~\exref{...}"
%   \extitle{slug}               the title from the docstring
%   \exerciseframe[opts]{slug}   beamer: the frames of the exercise
%   \studyexercise{slug}         study guide: heading, statement, solution
%   \scriptexercise{slug}        script: the statement box, as on the slide
%   \allexercises                exercises.tex: every exercise, in order; on
%                                the website statement and solution with
%                                markers where the pipeline splits the HTML
%                                (\collectionexercise), in the PDF plain:
%                                each statement, followed by its solution if
%                                \ifsolutions (exercises-solutions.tex)

\def\exweek{9}

\directlua{exercises = dofile("exercises/exercises.lua")
  exercises.scan("exercises", "\exweek")}

\ExplSyntaxOn

\msg_new:nnn { exercises } { unknown } { Unknown~exercise~`#1'. }

\cs_new_protected:Npn \__ex_missing:n #1
  { \msg_error:nnn { exercises } { unknown } {#1} ?? }
\cs_new:Npn \exref #1
  {
    \cs_if_exist:cTF { exnum@#1 }
      { \use:c { exnum@#1 } } { \__ex_missing:n {#1} }
  }
\cs_new:Npn \extitle #1
  {
    \cs_if_exist:cTF { extitle@#1 }
      { \use:c { extitle@#1 } } { \__ex_missing:n {#1} }
  }

% ---------------------------------------------------------------------------
% Loading an exercise: its parts are stored, nothing is typeset. The slides
% typeset the stored parts; the other documents read the part they need from
% the file again (\__ex_input_part:nn), as stored tables do not survive lwarp.

\tl_new:N \g__ex_current_tl
\tl_new:N \g__ex_problem_tl
\tl_new:N \g__ex_slideproblem_tl
\tl_new:N \g__ex_solution_tl
\seq_new:N \g__ex_slides_seq
\box_new:N \l__ex_scratch_box

\cs_new_protected:Npn \__ex_load:n #1
  {
    \cs_if_exist:cF { exfile@#1 } { \__ex_missing:n {#1} }
    \tl_gset:Nn \g__ex_current_tl {#1}
    \tl_gclear:N \g__ex_problem_tl
    \tl_gclear:N \g__ex_slideproblem_tl
    \tl_gclear:N \g__ex_solution_tl
    \seq_gclear:N \g__ex_slides_seq
    \vbox_set:Nn \l__ex_scratch_box
      {
        \DeclareDocumentEnvironment { problem } { +b }
          { \tl_gset:Nn \g__ex_problem_tl {##1} } { }
        \DeclareDocumentEnvironment { slideproblem } { +b }
          { \tl_gset:Nn \g__ex_slideproblem_tl {##1} } { }
        \DeclareDocumentEnvironment { slides } { +b }
          { \seq_gput_right:Nn \g__ex_slides_seq {##1} } { }
        \DeclareDocumentEnvironment { solution } { +b }
          { \tl_gset:Nn \g__ex_solution_tl {##1} } { }
        \exp_args:Nv \input { exfile@#1 }
      }
  }

% Typesets one part (problem, slideproblem, solution) of an exercise straight
% from its file; the other parts are skipped.
\cs_new_protected:Npn \__ex_input_part:nn #1#2
  {
    \group_begin:
      \clist_map_inline:nn { problem, slideproblem, slides, solution }
        {
          \str_if_eq:nnTF {##1} {#2}
            { \DeclareDocumentEnvironment {##1} { } { } { } }
            { \DeclareDocumentEnvironment {##1} { +b } { } { } }
        }
      \exp_args:Nv \input { exfile@#1 }
    \group_end:
  }

% Exercise 9.1, or Exercise 9.2 · Finance
\cs_new:Npn \__ex_heading:
  {
    Exercise~\exref { \g__ex_current_tl }
    \cs_if_exist:cT { exfinance@ \g__ex_current_tl } { ~·~Finance }
  }
\prg_new_conditional:Npnn \__ex_if_finance: { T, F, TF }
  {
    \cs_if_exist:cTF { exfinance@ \g__ex_current_tl }
      { \prg_return_true: } { \prg_return_false: }
  }

% ---------------------------------------------------------------------------
% Beamer

% The finance context is printed in the study guide, the exercise PDFs and on
% the website; the
% slides and the script (which shows what the slide shows) leave it out.
\bool_new:N \l__ex_context_bool
\bool_set_true:N \l__ex_context_bool
\DeclareDocumentEnvironment { context } { }
  {
    \bool_if:NTF \l__ex_context_bool
      { \par \textit { Finance~context. } ~ }
      { \__ex_gobble_context:w }
  }
  { \bool_if:NT \l__ex_context_bool { \par } }
\cs_new_protected:Npn \__ex_gobble_context:w #1 \end #2
  { \str_if_eq:nnTF {#2} { context } { \end { context } } { \__ex_gobble_context:w } }

\IfClassLoadedTF { beamer }
  {
    \bool_set_false:N \l__ex_context_bool

    \cs_new:Npn \__ex_block: { \__ex_if_finance:TF { exampleblock } { block } }
    \NewDocumentCommand \exstatement { }
      {
        \begin { \__ex_block: } { \__ex_heading: }
          \tl_if_empty:NTF \g__ex_slideproblem_tl
            { \g__ex_problem_tl } { \g__ex_slideproblem_tl }
        \end { \__ex_block: }
      }

    \tl_new:N \l__ex_frames_tl
    \NewDocumentCommand \exerciseframe { O{t} m }
      {
        \__ex_load:n {#2}
        \tl_clear:N \l__ex_frames_tl
        \seq_map_indexed_inline:Nn \g__ex_slides_seq
          {
            \tl_put_right:Nn \l__ex_frames_tl { \begin{frame}[#1] }
            \tl_put_right:Ne \l__ex_frames_tl
              {
                {
                  \exp_not:N \extitle {#2}
                  \int_compare:nNnT {##1} > { 1 } { \c_space_tl (cont.) }
                }
              }
            \tl_put_right:Nn \l__ex_frames_tl { ##2 \end{frame} }
          }
        \tl_use:N \l__ex_frames_tl
      }
  }
  {
    \cs_new:Npn \__ex_box: { \__ex_if_finance:TF { fin } { task } }
    \cs_new_protected:Npn \__ex_statement_box:n #1
      {
        \begin { \__ex_box: } { \__ex_heading: }
          \__ex_input_part:nn { \g__ex_current_tl } {#1}
        \end { \__ex_box: }
      }

    \NewDocumentCommand \studyexercise { m }
      {
        \__ex_load:n {#1}
        \subsection* { Exercise~\exref{#1}~---~\extitle{#1} }
        \addcontentsline { toc } { subsection }
          { Exercise~\exref{#1}~---~\extitle{#1} }
        \__ex_statement_box:n { problem }
        \soltitle \__ex_input_part:nn {#1} { solution }
      }

    \NewDocumentCommand \scriptexercise { m }
      {
        \__ex_load:n {#1}
        \group_begin:
          \bool_set_false:N \l__ex_context_bool
          \tl_if_empty:NTF \g__ex_slideproblem_tl
            { \__ex_statement_box:n { problem } }
            { \__ex_statement_box:n { slideproblem } }
        \group_end:
      }

    % The website (lwarp): statement and solution of each exercise, between
    % markers where the pipeline splits the HTML.
    \NewDocumentCommand \collectionexercise { m }
      {
        \__ex_load:n {#1}
        \par [[exercise:#1:problem]] \par
        \__ex_input_part:nn {#1} { problem }
        \par [[exercise:#1:solution]] \par
        \__ex_input_part:nn {#1} { solution }
        \par [[exercise:#1:end]] \par
      }
    % The PDF: plain, each statement, followed by its solution in the
    % version with solutions (\ifsolutions). The statement is set off like a
    % quote in Markdown, by a bar on the left (quote=<color>, also used by
    % exercises.tex) in the color of its box in the study guide: navy, or teal
    % for a finance exercise.
    \tcbset
      {
        quote/.style =
          {
            enhanced, breakable, frame~hidden, interior~hidden,
            sharp~corners, boxrule = 0pt, borderline~west = {2pt}{0pt}{#1},
            left = 8pt, right = 0pt, top = 1pt, bottom = 1pt,
            before~skip = 6pt, after~skip = 6pt,
            % the document's paragraph spacing (parskip=half)
            before~upper = { \setlength { \parskip } { .5\baselineskip } }
          }
      }
    \NewDocumentCommand \plainexercise { m }
      {
        \__ex_load:n {#1}
        \subsection* { Exercise~\exref{#1}~---~\extitle{#1} }
        \addcontentsline { toc } { subsection }
          { Exercise~\exref{#1}~---~\extitle{#1} }
        \__ex_if_finance:TF
          { \begin{tcolorbox}[quote=teal] }
          { \begin{tcolorbox}[quote=navy] }
          \__ex_input_part:nn {#1} { problem }
        \end{tcolorbox}
        \legacy_if:nT { solutions }
          {
            \soltitle
            \__ex_input_part:nn {#1} { solution }
          }
      }
    \NewDocumentCommand \allexercises { }
      {
        \IfPackageLoadedTF { lwarp }
          { \directlua { exercises.each("collectionexercise") } }
          { \directlua { exercises.each("plainexercise") } }
      }
  }

\ExplSyntaxOff


% PDF properties and the date on the title slide
% The properties of every PDF of Week 9: author (with email address), title,
% subject, keywords and the date the document last changed, in the PDF
% (hyperref) and on its title page (\date), and the header and footer of the
% printed documents (below). Each document loads this file after hyperref and
% then calls \documentinfo{<what it is>}, e.g. \documentinfo{Slides}.
%
% The date: the website build (scripts/build-content.mjs) puts
% \def\lastmodified{YYYY-MM-DD} in front of the document, the date of the last
% commit that changed its sources. Compiled locally, it is today.

\newcommand{\courseinfo}{Mathematics for Finance}
\newcommand{\weeklabel}{Week 9}
\newcommand{\weekinfo}{\weeklabel: Probability, conditional probability, independence, random variables}
\newcommand{\authorname}{Tibor SÁNDOR}
\providecommand{\myemail}{tiborsandor2001@gmail.com}

\makeatletter
\ifdefined\lastmodified\else
	\edef\lastmodified{\the\year-\two@digits{\the\month}-\two@digits{\the\day}}
\fi
\def\di@split#1-#2-#3\@nil{%
	\def\di@year{#1}\def\di@month{#2}\def\di@day{#3}}
\expandafter\di@split\lastmodified\@nil
% November 11, 2026
\edef\lastmodifiedtext{%
	\ifcase\numexpr\di@month\relax\or January\or February\or March\or April\or
		May\or June\or July\or August\or September\or October\or November\or
		December\fi\space\number\di@day, \di@year}
% D:20261111000000, the date format of PDF metadata
\edef\di@pdfdate{D:\di@year\di@month\di@day000000}

\newcommand{\documentinfo}[1]{%
	\hypersetup{
		pdftitle={\courseinfo, \weekinfo\space -- #1},
		pdfauthor={\authorname\space<\myemail>},
		pdfsubject={\courseinfo, \weekinfo},
		pdfkeywords={probability, combinatorics, mathematics for finance,
				Corvinus University of Budapest},
		pdfcreationdate={\di@pdfdate},
		pdfmoddate={\di@pdfdate},
		pdflang={en-US},
	}%
	\date{Last updated: \lastmodifiedtext}%
	\gdef\documentkind{#1}%
}

% Header and footer of the printed documents (KOMA-Script, scrlayer-scrpage;
% the slides have beamer's footline, the website neither): the course, the
% week and the document on the left of the header (the week's topic is on the
% title page), the university on the right; in the footer the author and the
% page as <current>/<total>. The title page gets only the footer.
\newcommand{\documentkind}{}
\IfClassLoadedTF{beamer}{}{\IfPackageLoadedTF{lwarp}{}{%
	\usepackage{lastpage}%
	\usepackage[headsepline=1.5pt,footsepline=1.5pt,plainfootsepline]
		{scrlayer-scrpage}%
	\clearpairofpagestyles
	\setkomafont{pageheadfoot}{\small\sffamily\color{navy}}%
	\setkomafont{pagenumber}{}%
	\setkomafont{headsepline}{\color{navy!35}}%
	\setkomafont{footsepline}{\color{navy!35}}%
	\ihead{\courseinfo\ \textperiodcentered\ \weeklabel\ \textperiodcentered\ 
		\documentkind}%
	\ohead{Corvinus University of Budapest}%
	\ifoot*{\authorname\ \textperiodcentered\ \myemail}%
	\ofoot*{\pagemark/\pageref*{LastPage}}%
	\pagestyle{scrheadings}%
	% Extra stretch for breaking lines with long inline formulas.
	\setlength{\emergencystretch}{1.5em}%
}}
\makeatother

\documentinfo{Slides}

\title[Mathematics for Finance]{Mathematics for Finance}
\author{Tibor SÁNDOR}
\subtitle{Week 9: Probability, conditional probability, independence,
  random variables}
\institute{Corvinus University of Budapest\\[2pt]
  \href{mailto:\myemail}{\texttt{\myemail}}}

\begin{document}

\begin{frame}
  \titlepage
\end{frame}

\begin{frame}{Plan for today}
  \begin{enumerate}
    \item \textbf{Session 1 (90 min):} probability of unions, intersections
          and complements; equally likely outcomes; conditional probability;
          Bayes' formula; odds.
    \item \textbf{Session 2 (90 min):} the generalized Bayes formula;
          independent events; discrete random variables, distribution
          function, expected value, variance and standard deviation.
  \end{enumerate}
  \vspace{.3cm}
  \textbf{Why does a finance student need this?} Default probabilities,
  early-warning models, betting odds, option payoffs and the risk of a
  strategy are all described with exactly these tools.
\end{frame}

%======================================================================
\section{Session 1: Probability and Bayes' formula}
%======================================================================

\begin{frame}
  \centering\Large\color{navy}Session 1: Probability, conditional
  probability, Bayes' formula
\end{frame}

\begin{frame}{From counting to probability}
  \small
  \begin{block}{Axioms}
    (1) $0\le P(E)\le1$; \quad (2) $P(\Om)=1$; \quad (3) if $E_1,E_2,\dots$
    are mutually exclusive, then $P\big(\bigcup_i E_i\big)=\sum_i P(E_i)$.
  \end{block}
  \begin{columns}[T]
    \begin{column}{.55\textwidth}
      \begin{block}{Rules that follow}
        $P(E^c)=1-P(E)$\\[2pt]
        $E\subseteq F\ \Rightarrow\ P(E)\le P(F)$\\[2pt]
        $P(EF^c)=P(E)-P(EF)$\\[2pt]
        $P(E\cup F)=P(E)+P(F)-P(EF)$
      \end{block}
    \end{column}
    \begin{column}{.43\textwidth}
      \centering
      \begin{tikzpicture}[scale=.62]
        \draw[thick] (-1.9,-1.45) rectangle (3.4,1.45);
        \node at (3.1,-1.15) {\scriptsize $\Om$};
        \draw[thick,navy] (0,0) circle (1.15);
        \draw[thick,teal] (1.5,0) circle (1.15);
        \node[navy] at (-1.55,1.05) {$E$};
        \node[teal] at (3.05,1.05) {$F$};
        \node at (-.45,0) {\scriptsize $EF^c$};
        \node at (.75,0) {\scriptsize $EF$};
        \node at (1.95,0) {\scriptsize $E^cF$};
        \node at (-1.35,-1.15) {\scriptsize $E^cF^c$};
      \end{tikzpicture}\\[1mm]
      \footnotesize the four regions are mutually exclusive;\\
      their probabilities add up to 1
    \end{column}
  \end{columns}
  \vspace{1mm}
  \textbf{Three events} (inclusion--exclusion, as for counts last week):
  \[ P(E\cup F\cup G)=P(E)+P(F)+P(G)-P(EF)-P(EG)-P(FG)+P(EFG). \]
\end{frame}

\exerciseframe{two-exams}

\begin{frame}{Equally likely outcomes}
  \begin{block}{Classical probability}
    If $\Om$ is finite and every outcome is equally likely, then
    \[ P(E)=\frac{|E|}{|\Om|}=\frac{\text{number of favorable outcomes}}
      {\text{number of all outcomes}}. \]
  \end{block}
  \begin{itemize}
    \item ``Equally likely'' must come from the setup: a fair die, a
          well-shuffled deck, a choice ``at random''.
    \item The counting methods of last week give $|E|$ and $|\Om|$.
    \item ``At least one'' $\to$ complement: $P(E)=1-P(E^c)$.
  \end{itemize}
\end{frame}

\exerciseframe{five-cards}

\begin{frame}{Conditional probability}
  \small
  \begin{columns}[T]
    \begin{column}{.58\textwidth}
      \begin{block}{Definition}
        If $P(F)>0$, the probability of $E$ \emph{given} $F$ is
        \[ P(E\mid F)=\frac{P(EF)}{P(F)}. \]
        $F$ becomes the new sample space; only the part $EF$ of $E$ counts.
      \end{block}
      \begin{block}{Multiplication rule}
        $P(EF)=P(F)\,P(E\mid F)$\\[2pt]
        $P(E_1E_2E_3)=P(E_1)\,P(E_2\mid E_1)\,P(E_3\mid E_1E_2)$
      \end{block}
    \end{column}
    \begin{column}{.4\textwidth}
      \centering
      \begin{tikzpicture}[scale=.62]
        \fill[teal!12] (1.5,0) circle (1.15);
        \begin{scope}
          \clip (0,0) circle (1.15);
          \fill[teal!45] (1.5,0) circle (1.15);
        \end{scope}
        \draw[thick] (-1.9,-1.45) rectangle (3.4,1.45);
        \node at (3.1,-1.15) {\scriptsize $\Om$};
        \draw[thick,navy] (0,0) circle (1.15);
        \draw[thick,teal] (1.5,0) circle (1.15);
        \node[navy] at (-1.55,1.05) {$E$};
        \node[teal] at (3.05,1.05) {$F$};
        \node at (.75,0) {\scriptsize $EF$};
      \end{tikzpicture}\\[1mm]
      \footnotesize $P(E\mid F)$: the share of $EF$ within $F$
    \end{column}
  \end{columns}
  \vspace{1mm}
  $P(\,\cdot\mid F)$ is a probability itself, e.g.\
  $P(E^c\mid F)=1-P(E\mid F)$.
\end{frame}

\exerciseframe{exams-conditional}

\exerciseframe{actuarial-exams}

\begin{frame}{Total probability and Bayes' formula}
  \small
  \begin{columns}[T]
    \begin{column}{.56\textwidth}
      \begin{block}{Law of total probability}
        \[ P(E)=P(E\mid F)\,P(F)+P(E\mid F^c)\,P(F^c) \]
      \end{block}
      \begin{block}{Bayes' formula}
        \[ P(F\mid E)=\frac{P(E\mid F)\,P(F)}
          {P(E\mid F)\,P(F)+P(E\mid F^c)\,P(F^c)} \]
      \end{block}
      We know $P(E\mid F)$, but we need $P(F\mid E)$: Bayes' formula turns
      the condition around.
    \end{column}
    \begin{column}{.42\textwidth}
      \centering
      \begin{tikzpicture}[every node/.style={font=\small},
          lab/.style={font=\scriptsize,inner sep=1pt,sloped},
          grow=right,level distance=2.3cm,
          level 1/.style={sibling distance=2.5cm},
          level 2/.style={sibling distance=1.25cm}]
        \node[circle,fill=navy,inner sep=2pt]{}
          child{node[fill=navy!10]{$F^c$}
            child{node{$E^c$} edge from parent node[lab,below]{$P(E^c|F^c)$}}
            child{node[fill=gold!30]{$E$} edge from parent node[lab,above]{$P(E|F^c)$}}
            edge from parent node[lab,below]{$P(F^c)$}}
          child{node[fill=navy!10]{$F$}
            child{node{$E^c$} edge from parent node[lab,below]{$P(E^c|F)$}}
            child{node[fill=gold!30]{$E$} edge from parent node[lab,above]{$P(E|F)$}}
            edge from parent node[lab,above]{$P(F)$}};
      \end{tikzpicture}\\[1mm]
      \footnotesize multiply along a path;\\
      add the paths that end in $E$
    \end{column}
  \end{columns}
\end{frame}

\exerciseframe{early-warning}

\begin{frame}{Odds}
  \begin{block}{Odds of an event}
    \[ \text{odds of } H=\frac{P(H)}{P(H^c)}=\frac{P(H)}{1-P(H)},
      \qquad P(H)=\frac{\text{odds}}{1+\text{odds}}. \]
    $P(H)=2/3$: odds $2$, ``2 to 1 in favor''. \quad $P(H)=0.2$: odds
    $1/4$, ``4 to 1 against''.
  \end{block}
  \begin{block}{Bayes' formula in odds form}
    \[ \underbrace{\frac{P(H\mid E)}{P(H^c\mid E)}}_{\text{new odds}}
      =\underbrace{\frac{P(H)}{P(H^c)}}_{\text{old odds}}\cdot
      \underbrace{\frac{P(E\mid H)}{P(E\mid H^c)}}_{\text{likelihood ratio}} \]
  \end{block}
  Evidence that is more likely if $H$ is true raises the odds of $H$.
\end{frame}

\exerciseframe{early-warning-odds}

\begin{frame}{Session 1: summary}
  \small
  \begin{itemize}
    \item $P(E^c)=1-P(E)$; \ $P(E\cup F)=P(E)+P(F)-P(EF)$;
          \ $P(EF^c)=P(E)-P(EF)$.
    \item Equally likely outcomes: $P(E)=|E|/|\Om|$; ``at least one''
          $\to$ complement.
    \item $P(E\mid F)=P(EF)/P(F)$; multiplication rule: multiply along a
          path of a tree.
    \item Total probability: add the paths that end in $E$. Bayes' formula:
          one path divided by all of them.
    \item Odds: new odds $=$ old odds $\times$ likelihood ratio.
    \item Rare events and imperfect signals: most alarms are false.
  \end{itemize}
\end{frame}

%======================================================================
\section{Session 2: Independence, random variables}
%======================================================================

\begin{frame}
  \centering\Large\color{navy}Session 2: Independence and discrete random
  variables
\end{frame}

\begin{frame}{Recap of Session 1}
  \begin{itemize}
    \item $P(E\cup F)=P(E)+P(F)-P(EF)$, \ $P(E^c)=1-P(E)$.
    \item $P(E\mid F)=\dfrac{P(EF)}{P(F)}$, \ $P(EF)=P(F)\,P(E\mid F)$.
    \item Two hypotheses $F$, $F^c$: total probability and Bayes' formula,
          best drawn as a tree.
    \item Odds: new odds $=$ old odds $\times$ likelihood ratio.
  \end{itemize}
  \vspace{2mm}
  \textbf{Now:} more than two hypotheses, independent events, and random
  variables: their distribution, expected value and standard deviation.
\end{frame}

\begin{frame}{The generalized Bayes formula}
  \small
  \begin{columns}[T]
    \begin{column}{.6\textwidth}
      $F_1,\dots,F_n$ is a \textbf{partition} of $\Om$: mutually exclusive
      events, and exactly one of them occurs.
      \begin{block}{Total probability}
        \[ P(E)=\sum_{i=1}^n P(E\mid F_i)\,P(F_i) \]
      \end{block}
      \begin{block}{Generalized Bayes formula}
        \[ P(F_j\mid E)=\frac{P(E\mid F_j)\,P(F_j)}
          {\sum_{i=1}^n P(E\mid F_i)\,P(F_i)} \]
      \end{block}
    \end{column}
    \begin{column}{.38\textwidth}
      \centering
      \begin{tikzpicture}[every node/.style={font=\small},
          lab/.style={font=\scriptsize,inner sep=1pt,sloped},
          grow=right,level distance=1.8cm,
          level 1/.style={sibling distance=1.45cm}]
        \node[circle,fill=navy,inner sep=2pt]{}
          child{node[fill=navy!10]{$F_3$}
            child[level distance=1.9cm]{node[fill=gold!30]{$E$}
              edge from parent node[lab,above]{$P(E|F_3)$}}
            edge from parent node[lab,below]{$P(F_3)$}}
          child{node[fill=navy!10]{$F_2$}
            child[level distance=1.9cm]{node[fill=gold!30]{$E$}
              edge from parent node[lab,above]{$P(E|F_2)$}}
            edge from parent node[lab,above]{$P(F_2)$}}
          child{node[fill=navy!10]{$F_1$}
            child[level distance=1.9cm]{node[fill=gold!30]{$E$}
              edge from parent node[lab,above]{$P(E|F_1)$}}
            edge from parent node[lab,above]{$P(F_1)$}};
      \end{tikzpicture}\\[1mm]
      \footnotesize one path to $E$ for each $F_i$
    \end{column}
  \end{columns}
\end{frame}

\exerciseframe{default-by-rating}

\begin{frame}{Independent events}
  \small
  \begin{block}{Two events}
    $E$ and $F$ are \textbf{independent} if $P(EF)=P(E)\,P(F)$. If
    $P(F)>0$, this means $P(E\mid F)=P(E)$: knowing $F$ does not change the
    chances of $E$. Then $E$ and $F^c$ are independent too.
  \end{block}
  \begin{block}{Three events}
    $E$, $F$, $G$ are independent if all four product rules hold:
    \[ P(EF)=P(E)P(F),\ \ P(EG)=P(E)P(G),\ \ P(FG)=P(F)P(G),\ \
      P(EFG)=P(E)P(F)P(G). \]
  \end{block}
  \begin{itemize}
    \item Mutually exclusive events with positive probabilities are
          \emph{never} independent: if one occurs, the other cannot.
    \item $n$ independent trials, each a success with probability $p$:
          \ $P(\text{at least one success})=1-(1-p)^n$.
  \end{itemize}
\end{frame}

\exerciseframe{pairwise-not-mutual}

\exerciseframe{bond-defaults}

\begin{frame}{Discrete random variables}
  \small
  A \textbf{random variable} $X$ assigns a number to every outcome. It is
  \textbf{discrete} if it takes finitely or countably many values.
  \begin{columns}[T]
    \begin{column}{.52\textwidth}
      \begin{block}{Probability mass function}
        $p(x)=P(X=x)$; \ $p(x)\ge0$, \ $\sum_x p(x)=1$.
      \end{block}
      \begin{block}{Distribution function}
        $F(x)=P(X\le x)=\sum_{y\le x}p(y)$.\\[2pt]
        A step function: it jumps by $p(x)$ at every value $x$ and is
        constant in between; it rises from 0 to 1.
      \end{block}
    \end{column}
    \begin{column}{.46\textwidth}
      \centering
      \begin{tikzpicture}[xscale=.85,yscale=2.1]
        \draw[->] (-.3,0)--(4.6,0) node[right,font=\tiny]{$x$};
        \draw[->] (0,-.05)--(0,1.2) node[above,font=\tiny]{$F(x)$};
        \foreach \y/\l in {.25/{1/4},.75/{3/4},.875/{7/8},1/1}
          \draw (-.05,\y)--(.05,\y) node[left=2pt,font=\tiny,fill=white,inner sep=.5pt] {$\l$};
        \foreach \x/\a/\b in {1/0/.25,2/.25/.75,3/.75/.875,4/.875/1}
          \draw[gray,densely dotted] (\x,\a)--(\x,\b);
        \draw[thick,teal] (-.3,0)--(1,0);
        \draw[thick,teal] (1,.25)--(2,.25) (2,.75)--(3,.75)
          (3,.875)--(4,.875) (4,1)--(4.5,1);
        \foreach \x/\y in {1/.25,2/.75,3/.875,4/1}
          \fill[teal] (\x,\y) circle (1.6pt);
        \foreach \x/\y in {1/0,2/.25,3/.75,4/.875}
          \draw[teal,fill=white] (\x,\y) circle (1.6pt);
        \foreach \x in {1,2,3,4} \node[below,font=\tiny] at (\x,0) {\x};
      \end{tikzpicture}\\[1mm]
      \footnotesize $p(1)=\frac14$, $p(2)=\frac12$, $p(3)=p(4)=\frac18$
    \end{column}
  \end{columns}
  \vspace{1mm}
  $P(a<X\le b)=F(b)-F(a)$, \quad $P(X>a)=1-F(a)$.
\end{frame}

\exerciseframe{urn-game}

\begin{frame}{Expected value, variance, standard deviation}
  \small
  \begin{block}{Expected value}
    $\E[X]=\sum_x x\,p(x)$: the average of the values, weighted by their
    probabilities.\\[2pt]
    $\E[g(X)]=\sum_x g(x)\,p(x)$, \qquad $\E[aX+b]=a\,\E[X]+b$.
  \end{block}
  \begin{block}{Variance and standard deviation}
    With $\mu=\E[X]$:
    \[ \Var(X)=\E\big[(X-\mu)^2\big]=\E[X^2]-\big(\E[X]\big)^2,\qquad
      \SD(X)=\sqrt{\Var(X)}. \]
    $\Var(aX+b)=a^2\Var(X)$, \qquad $\SD(aX+b)=|a|\,\SD(X)$.
  \end{block}
  In finance: $\E[X]$ is the expected return or payoff, $\SD(X)$ measures
  the \emph{risk}.
\end{frame}

\exerciseframe{call-option}

\exerciseframe{trading-strategy}

\begin{frame}{Session 2: summary}
  \small
  \begin{itemize}
    \item Partition $F_1,\dots,F_n$: $P(E)=\sum_i P(E\mid F_i)P(F_i)$, and
          the generalized Bayes formula gives $P(F_j\mid E)$.
    \item Independence: $P(EF)=P(E)P(F)$; for three events all four
          product rules; pairwise independence is not enough.
    \item pmf $p(x)=P(X=x)$; distribution function $F(x)=P(X\le x)$, a step
          function.
    \item $\E[X]=\sum x\,p(x)$; \ $\E[g(X)]=\sum g(x)p(x)$, in general not
          $g(\E[X])$.
    \item $\Var(X)=\E[X^2]-(\E[X])^2$; \ $\Var(aX+b)=a^2\Var(X)$; \
          $\SD(X)=\sqrt{\Var(X)}$.
  \end{itemize}
\end{frame}

\end{document}
